% \section{Discussion}\label{sec:discussion}

% This study identifies an everyday interactional route through which human--LLM use can become a setting for informal learning. Learning-oriented engagement was common across ordinary coding and writing conversations, but its constructive form was selective: it concentrated in explicit learning intent, coding tasks and deeper exchanges. Assistant scaffolding marked conversations with richer constructive participation, yet the association depended on support form and temporal scale. Feedback-like and explanatory moves were more consistently linked to constructive engagement than directive or questioning moves, and the strongest temporal evidence appeared locally, from one turn to the next, rather than as a simple cumulative rise across a conversation. Together, these findings show that informal learning in everyday LLM use is not a generic consequence of exposure to AI; it emerges when users and assistants sustain iterative work with ideas.

% The main finding is therefore not simply that engagement exists, but that constructive engagement is ecological. It is amplified by explicit learning intent, strongly structured by task domain and concentrated in deeper interactional sequences. This depth pattern should be read as an opportunity marker: longer conversations provide more occasions for constructive turns and may reflect more complex tasks or more engaged users, but they do not by themselves imply a cumulative learning trajectory. Coding produced substantially more cognitive and constructive engagement than writing, even when conversations were not explicitly framed as learning. This suggests that task ecology can shape the baseline opportunity for informal learning as much as users' stated intentions do. For research on human--LLM interaction, the implication is that everyday LLM use should not be treated as a single behavioural environment. Different domains organize different opportunities for users to ask, test, refine and extend ideas.

% Assistant scaffolding was associated with richer constructive participation, but the association was heterogeneous rather than treatment-like. Conversations containing scaffolded support had higher constructive ratios and greater post-answer depth across all four matched settings, and adjusted models showed that these associations were not explained away by observed covariates. In coding, scaffolded conversations also showed higher persistence after difficulty. These results suggest that scaffolding is part of the interactional environment in which constructive participation becomes more likely. At the same time, the observational design prevents causal claims about instructional efficacy. The contribution is to identify structured associations in naturalistic use, not to show that exposure to scaffolding produces durable learning gains.

% A central implication is that increasing engagement is not a matter of increasing scaffolding in general. Support form mattered. Feedback-like and explanatory moves were the most consistently positive forms for constructive engagement, whereas instructing and questioning were weaker or negative. Support supply also differed by domain: coding conversations were explanation-heavy, while writing conversations were more directive. This two-sided heterogeneity helps explain why support did not have a single portable engagement value. The same broad category of scaffolding can invite different forms of user participation depending on how it is delivered and on the task ecology in which it appears.

% The temporal analyses further suggest that informal learning with LLMs is best understood as a local interaction process. Scaffolded turns were followed by higher constructive uptake in most settings, and users' current engagement state shaped the probability of receiving subsequent scaffolding. Support and engagement therefore appeared as coupled states, not as independent events in a one-way instructional pipeline. This temporal pattern also guards against over-reading the conversation-level associations: the data show local reciprocal ordering, not experimentally identified amplification. Longer-span trajectories were more mixed: around-first-scaffold contrasts were small or negative in several settings, whereas first-half versus second-half analyses in WildChat showed increases in constructive engagement. The strongest temporal claim is therefore scale-specific. Everyday LLM interaction repeatedly creates local opportunities for constructive uptake, but those opportunities do not automatically accumulate into a uniform upward trajectory.

% These findings clarify why everyday LLMs are neither merely tools for task completion nor full substitutes for designed instruction. In ordinary use, the same system may act as an assistant, explainer, feedback provider or directive instructor within a single conversation. The learning relevance of the interaction depends on whether these roles recruit further user thinking rather than simply resolving the immediate task. This perspective also explains why constructive engagement is a useful organizing outcome. It captures the moments when users are not only receiving information but actively working with it: testing a boundary, revising an idea, articulating a constraint or extending an explanation.

% For design, the results point toward adaptive support that is calibrated to user state, task ecology and support form. Systems that aim to support informal learning should not simply add more instructional language. They should identify when a user is already constructively engaged and sustain that trajectory, when a user is active but not yet constructive and offer feedback or explanation that invites elaboration, and when a user primarily needs direct help without forcing a pedagogical detour. The domain contrast is important here. Coding often benefited from explanatory support and showed stronger persistence effects, whereas writing contained more directive support and weaker temporal gains. Future systems may therefore need domain-sensitive support policies rather than a single generic tutor mode.

% Several limits define the boundary of the claims. First, the analyses are observational, so scaffolding should be interpreted as associated with engagement rather than causally responsible for it. Second, the labels capture behavioural manifestations of learning-oriented engagement rather than verified learning outcomes or durable knowledge change. Third, fine-grained support categories are noisier than the top-level distinction between direct response and scaffolding, and means labels are not mutually exclusive. Fourth, WildChat model labels provide useful robustness checks but are not experimental assignments; model comparisons may be confounded with user population, time and task mix. These limits are important, but they do not erase the central pattern: everyday human--LLM interaction contains structured, context-dependent dynamics characteristic of informal learning.

% The most direct unresolved threat is selection into support: more complex tasks or more motivated users may both elicit scaffolding and produce constructive follow-up. The adjacent-turn analyses reduce this concern by conditioning on the immediately preceding interaction state, but they do not eliminate unobserved time-varying confounding. The safest interpretation is therefore local, state-dependent coupling between assistant support and user engagement rather than cumulative uplift or causal learning gain.

% Overall, the study identifies a middle space in contemporary AI use. LLMs are not only instruments for faster task completion, and they are not only educational technologies when explicitly deployed as tutors. They are increasingly part of the ordinary interactional settings in which people clarify, test, revise and extend understanding while working. Studying this middle space is necessary for designing future systems that support learning without assuming that every user has entered a classroom, and for understanding how general-purpose AI may reshape informal learning at scale.


\section{Discussion}

% This study identifies everyday human--LLM interaction as a middle space between ordinary tool use and designed instruction. Across naturalistic coding and writing conversations, learning-oriented engagement was recurrent, but its constructive form was selective. Users often used, applied or followed up on assistant responses, yet deeper constructive engagement emerged when the exchange gave users reasons and opportunities to keep working with ideas---through explicit learning framing, task affordances that made testing and revision salient, and assistant support that left room for user evaluation rather than simply completing the task. Assistant scaffolding was therefore not informative as a uniform support exposure. It was associated with richer constructive participation when its form, timing and surrounding user state supported continued sense-making: feedback and explanation aligned more clearly with constructive engagement than more directive or question-based moves, and temporal evidence was strongest in adjacent-turn coupling rather than cumulative uplift. Together, these findings suggest that informal learning in everyday LLM use is a situated interactional achievement, not a generic consequence of AI access: it emerges when users and assistants maintain a shared space for testing, revising and extending understanding while pursuing practical goals.



% This study identifies everyday human--LLM interaction as a middle space between ordinary tool use and designed instruction. In naturalistic coding and writing conversations, learning-oriented engagement was recurrent, but its constructive form was selective. Users often used, applied or followed up on assistant responses; deeper constructive engagement emerged when the exchange preserved reasons and opportunities for users to keep working with ideas---through explicit learning framing, task affordances that made testing and revision salient, and support that left room for user evaluation rather than simply completing the task. Assistant scaffolding was therefore not informative as a uniform support exposure. It was associated with richer constructive participation when its form, timing and surrounding user state supported continued sense-making: feedback and explanation aligned more clearly with constructive engagement than more directive or question-based moves, and temporal evidence was strongest in adjacent-turn coupling rather than cumulative uplift. Together, these findings suggest that informal learning in everyday LLM use is a situated interactional achievement, not a generic consequence of AI access: it emerges when users and assistants maintain a shared space for testing, revising and extending understanding while pursuing practical goals.




% comment: 感觉这一段的开头insights里关于scaffolding的部分还是有一点点浅？"and when the interaction extended beyond one-off answer delivery" 似乎并没有和前面的"Users often took up, applied or followed up on assistant responses" 有太大的区别，是不是得从assistant的scaffolding的角度说一点一句话的但又很insightful的内容?然后和下面的assistant scaffolding接到一起


This study identifies everyday human--LLM interaction as a middle space between ordinary tool use and designed instruction. In naturalistic coding and writing conversations, learning-oriented engagement was recurrent, but its constructive form was selective. Users often used, applied or followed up on assistant responses; deeper constructive engagement emerged when the exchange preserved reasons and opportunities for users to keep working with ideas---through explicit learning framing, task affordances that made testing and revision salient, and support that left room for user evaluation rather than simply completing the task. Assistant scaffolding was therefore not informative as a uniform support exposure. It was associated with richer constructive participation when its form, timing and surrounding user state supported continued sense-making: feedback and explanation aligned more clearly with constructive engagement than more directive or question-based moves, and temporal evidence was strongest in adjacent-turn coupling rather than cumulative uplift. Together, these findings suggest that informal learning in everyday LLM use is a situated interactional achievement, not a generic consequence of AI access: it emerges when users and assistants maintain a shared space for testing, revising and extending understanding while pursuing practical goals.

The results suggest that constructive engagement in everyday AI use is ecological. It is not simply a property of an individual user, nor a direct consequence of a model producing educational content. Instead, it depends on how the user's goal, the task ecology and the unfolding exchange configure opportunities for sense-making. Explicit learning framing amplified constructive engagement, but constructive participation also appeared in task-oriented conversations where learning was incidental. This matters because informal learning rarely announces itself as a formal learning episode. It often appears when a practical task exposes uncertainty, invites comparison or requires revision. Coding and writing illustrate this point as task ecologies rather than as isolated topics. Some tasks make breakdowns, constraints and tests more visible, while others allow users to accept or apply a revision without articulating the rationale. Similarly, sustained exchanges should not be treated as evidence of learning progression by themselves; they are opportunity structures in which users can return to, interrogate and reshape the assistant's response. Everyday LLM use should therefore not be analysed as a single behavioural environment. Different tasks organize different possibilities for users to ask, test, refine and extend ideas.

The support results further suggest that scaffolding is best understood as an interactional condition rather than a treatment-like exposure. Conversations containing scaffolded support showed richer constructive participation, greater post-answer depth and, in coding-oriented conversations, greater persistence after observable breakdowns. These patterns indicate that scaffolding is part of the conversational environment in which constructive participation becomes more likely to appear. At the same time, scaffolded support was not randomly assigned. More difficult tasks, more motivated users or already constructive exchanges may elicit scaffolding while also producing constructive follow-up. The point is therefore not that exposure to scaffolding causes durable learning gains. Rather, the study identifies structured associations in naturalistic interaction and shows where they become visible: in conversations where support and user participation are mutually organized around continued work with ideas.

The most actionable insight is that learning-oriented support is not a matter of adding more scaffolding in general. The broad label of scaffolding concealed important differences in support form. Feedback-like support and explanation were most clearly aligned with constructive engagement, whereas instructing and questioning were less aligned with visible elaboration. This distinction suggests that different support forms leave different roles for the user. Feedback responds to a user's prior attempt and can create an occasion for revision, evaluation or extension. Explanation can make mechanisms visible and give users material for testing their understanding. By contrast, detailed instructions can make the next step executable without requiring the user to articulate why it works, and questions in everyday assistant use may function as information gathering or clarification requests rather than carefully designed prompts for reflection. This does not mean that instruction or questioning are unhelpful. They may be appropriate when users need direct task support, when information is missing or when the cost of continued exploration is high. But they illustrate why support quality for informal learning cannot be judged only by whether a response is pedagogical in form. The relevant question is whether the support preserves a productive role for the user in the work of understanding.

% The temporal results clarify how this form of informal learning differs from, and can complement, formal instruction. Everyday LLM interactions repeatedly create local moments in which support and user engagement align, but those moments do not automatically accumulate into a coherent learning trajectory. This is a central distinction between informal AI-mediated learning and formal learning environments. Formal instruction can provide curricular sequencing, shared standards, assessment, accountability and social structures that help local learning episodes accumulate over time. Everyday AI use instead operates at the point of need: it enters moments of confusion, repair, comparison and revision while users are trying to solve real problems. Its value is therefore not to replace formal education, but to occupy the spaces between formal learning episodes, where people encounter unfamiliar concepts and decide what to try next. The temporal findings show both the promise and the limitation of this role. LLM interaction can create situated opportunities for constructive engagement, but those opportunities require additional design support if they are to become durable knowledge, transferable skill or cumulative expertise.

The temporal results clarify a division of labour between AI-mediated informal learning and formal instruction. Everyday LLM interactions repeatedly create local moments in which support and user engagement align, but those moments do not automatically accumulate into a coherent learning trajectory. This pattern does not mean that everyday AI use should be judged by whether it reproduces formal instruction. Rather, it points to a complementary role. Formal instruction is designed to accumulate learning over time through curricular sequencing, shared standards, assessment, accountability and opportunities to abstract beyond the immediate task. Everyday AI use is more episodic and situated: it enters moments when users are confused, blocked, comparing alternatives, repairing a failed attempt or adapting knowledge to a concrete problem. Its value is therefore not to replace formal education, but to extend learning into the spaces between and around formal learning episodes, where people encounter unfamiliar concepts while doing real work. The temporal findings show both the promise and the limitation of this role. LLM interaction can create situated opportunities for constructive engagement, but those opportunities require additional design support if they are to become durable knowledge, transferable skill or cumulative expertise. Future assistants may therefore need to bridge informal and formal learning functions: preserving the immediacy and personalization of point-of-need support, while helping users connect repeated local episodes into longer arcs of reflection, consolidation and skill development.

% The temporal results clarify how AI-mediated informal learning differs from, and can complement, formal instruction. Everyday LLM interactions repeatedly create local moments in which support and user engagement align, but those moments do not automatically accumulate into a coherent learning trajectory. This is not a weakness of informal learning alone; it reflects a different role in the learning ecosystem. Formal instruction provides what local interaction often lacks: curricular sequencing, shared standards, assessment, accountability and opportunities to abstract beyond the immediate problem. Everyday AI use, by contrast, operates at the point of need. It enters moments when users are confused, blocked, comparing alternatives, repairing a failed attempt or trying to adapt knowledge to a concrete task. Its value is therefore not to replace formal education, but to extend learning into the spaces where formal instruction is absent: the moments when people encounter unfamiliar concepts while doing real work. The temporal findings show both the promise and the limitation of this role. LLM interaction can create situated opportunities for constructive engagement, but those opportunities require additional design support if they are to become durable knowledge, transferable skill or cumulative expertise. Future assistants may therefore need to bridge informal and formal learning functions: preserving the immediacy and personalization of point-of-need support, while helping users connect repeated local episodes into longer arcs of reflection, consolidation and skill development.

% These findings have implications for personal AI assistants that aim to support lifelong learning in the flow of everyday work. The design target is not a generic tutor mode, and it is not simply more instructional language. A useful assistant must decide what kind of role to leave for the user at a given moment. Sometimes the right action is to answer directly, because the user primarily needs task completion or because a pedagogical detour would impose unnecessary friction. Sometimes the right action is to explain a mechanism, because the user is trying to understand why something works. Sometimes it is to give feedback on a user's attempt, because the user has already produced an idea that can be evaluated, revised or extended. Sometimes it is to offer a hint or partial cue, preserving the user's responsibility for the next step. And sometimes it may be to step back, because additional instruction would replace rather than support the user's own sense-making. More ambitiously, assistants could help connect local learning moments over time: noticing recurring misconceptions, revisiting earlier explanations, prompting reflection after repeated failures or helping users consolidate what they learned from a practical task. This would move AI support beyond answer delivery and beyond one-off tutoring moves, toward lifelong learning assistance that helps users build knowledge while solving real problems.

% These findings suggest a design principle for personal AI assistants that aim to support lifelong learning: support should be calibrated by the cognitive role left to the user, not by a generic tutor mode. The relevant design question is not simply whether to provide help, but whether the current exchange calls for answer delivery, explanation, feedback, a partial cue or a deliberate step back. When users mainly need completion, direct help may be appropriate; when they reveal uncertainty or offer an attempted solution, feedback and explanation can keep the repair process with the user by making mechanisms and evaluative criteria visible; when users are close to a solution, hints may preserve agency without replacing their reasoning. Questions should likewise be used purposefully rather than as a default tutoring script, because in everyday assistant use they may impose conversational burden rather than reflection. The larger challenge is longitudinal. Because local support--engagement alignments do not automatically accumulate, future assistants should help users connect repeated episodes by tracking recurring misconceptions, revisiting prior explanations and prompting consolidation after difficult tasks. Such systems would support lifelong learning not by turning every interaction into a lesson, but by converting moments of problem solving into opportunities for durable understanding.

% This complementarity suggests a design principle for personal AI assistants that support lifelong learning: they should not simply imitate a tutor, but should connect situated problem solving with longer arcs of understanding. At the moment-to-moment scale, the key design question is how much cognitive responsibility should remain with the user. Direct answers may be appropriate when the user primarily needs task completion; explanations may be appropriate when the user is trying to understand a mechanism; feedback may be most useful when the user has produced an attempt that can be evaluated and revised; hints may preserve agency when the user is close to a solution; and stepping back may be necessary when additional instruction would replace the user's own sense-making. At the longitudinal scale, however, the assistant must do more than choose the right support form in a single exchange. Because local support--engagement alignments do not automatically accumulate, future systems should help users connect repeated episodes across time: recognizing recurring misconceptions, revisiting prior explanations, prompting reflection after difficult tasks and helping users consolidate what a practical episode revealed. In this sense, the promise of personal AI assistants is not to turn every interaction into formal instruction, but to provide continuity across the informal learning moments that occur while people solve real problems.

This complementarity suggests a design principle for personal AI assistants that support lifelong learning: they should not simply imitate formal tutors or add instructional language, but should manage the cognitive responsibility left to the user across time. At the local scale, support should be calibrated to the user's state and the task affordance: direct help may be appropriate when completion is the immediate goal, whereas explanation, feedback or partial cues are more useful when the user is trying to diagnose, evaluate or revise an idea. Equally important is restraint, because too much instruction can replace the user's own sense-making rather than support it. At the longitudinal scale, the assistant must do more than choose the right support form in a single exchange. Because local support--engagement alignments do not automatically accumulate, future systems should help users connect repeated episodes across time: recognizing recurring uncertainties, linking current problems to prior explanations, and turning repeated breakdowns into occasions for reflection and consolidation. In this sense, the distinctive promise of lifelong AI assistance is not to convert everyday work into formal instruction, but to give continuity, memory and personalization to the informal learning that happens between formal learning episodes. Together with formal instruction, such systems could support a more continuous and personalized learning ecosystem in which schools and curricula provide structure, standards and shared progression, while personal AI assistants carry learning into everyday practice, helping people transform formal knowledge into situated judgment, transferable skill and cumulative understanding.


% Several limits define the boundary of these claims. First, the analyses are observational, so scaffolded support should be interpreted as associated with engagement rather than causally responsible for it. Selection into support remains a central threat: complex tasks or motivated users may both elicit scaffolding and produce constructive follow-up. The adjacent-turn analyses reduce some ambiguity by conditioning on the immediately preceding interaction state, but they cannot eliminate unobserved time-varying confounding. Second, the labels capture behavioural manifestations of learning-oriented engagement rather than verified learning outcomes or durable knowledge change. Constructive engagement is a theoretically grounded behavioural proxy, not a direct measure of what users learned. Third, fine-grained support-form labels are not mutually exclusive and are noisier than the top-level distinction between scaffolded and non-scaffolded support. Fourth, the analysis focuses on coding and writing as two common everyday task ecologies; other domains may organize different opportunities for engagement and support. Finally, platform and model labels are not experimental assignments, so platform- or model-specific patterns should be treated as robustness evidence rather than causal comparisons. These limits define the appropriate interpretation: the study does not show that LLM support causes learning, but it shows how learning-relevant interaction can be observed, characterized and designed for in everyday AI use.

Several boundaries define the scope of these claims. The analyses use naturalistic logs, so they identify structured associations rather than causal effects. Scaffolded support was not assigned, and conversations that elicited support may also have involved more complex tasks, more motivated users or prior engagement. The data also capture interactional episodes rather than long-term learning histories: they do not show whether repeated use leads individual users to retain knowledge, transfer skills or become more independent over time. In addition, engagement labels are behavioural proxies. Constructive engagement marks observable elaboration, testing, revision or extension, but it is not a direct measure of durable learning outcomes. Finally, the corpus covers large public WildChat sample covering two common task ecologies, coding and writing; other domains, populations and deployment contexts may organize different opportunities for engagement and support. Future work should therefore combine large-scale log analysis with controlled interventions, longitudinal user-level designs and independent measures of knowledge gain or skill transfer.

% Overall, the study suggests that LLMs are becoming part of the ordinary environments in which people encounter, test and extend understanding while working. They are not only instruments for faster task completion, and they are not educational technologies only when explicitly deployed as tutors. They are increasingly embedded in the everyday sequences through which people clarify problems, revise ideas and decide what to try next. Understanding this middle space is necessary for building future assistants that support lifelong learning without assuming that every user has entered a classroom. The central challenge is to make everyday human--AI collaboration a setting in which people can solve real problems while continuing to build, test and extend what they understand.

Within these boundaries, the study makes everyday AI-supported informal learning empirically visible. It does not show that LLM support causes learning; instead, it shows where learning-oriented engagement appears, how its constructive form is shaped by task ecology and user framing, and why assistant support must be understood through form, timing and local user state rather than through scaffolding presence alone. This reframes everyday human--LLM interaction as a middle space between ordinary tool use and formal instruction. Together with formal instruction, such systems could support a more continuous and personalized learning ecosystem in which curricula provide structure, standards and shared progression, while personal AI assistants extend learning into everyday practice, helping learners transform situated problem solving into durable understanding, transferable expertise and lifelong learning agency.

%%%%%%%%%%%%%%%%%%%%%%%%%%%%%%%%%%%%%%%%%%%%%%%%%%%%%


%不要直接表述成contribution这样是不是更好

% A first contribution is to show that constructive engagement in everyday AI use is ecological. It is not simply a property of an individual user, nor a direct consequence of a model producing educational content. Instead, it depends on how the user's goal, the task ecology and the unfolding exchange configure opportunities for sense-making. Explicit learning framing amplified constructive engagement, but constructive participation also appeared in task-oriented conversations where learning was incidental. This matters because informal learning rarely announces itself as a formal learning episode. It often appears when a practical task exposes uncertainty, invites comparison or requires revision. Coding and writing illustrate this point as task ecologies rather than as isolated topics. Some tasks make breakdowns, constraints and tests more visible, while others allow users to accept or apply a revision without articulating the rationale. Similarly, sustained exchanges should not be treated as evidence of learning progression by themselves; they are opportunity structures in which users can return to, interrogate and reshape the assistant's response. The broader implication is that everyday LLM use should not be analysed as a single behavioural environment. Different tasks organize different possibilities for users to ask, test, refine and extend ideas.

% A second contribution is to reframe scaffolding as an interactional condition rather than a treatment-like exposure. Conversations containing scaffolded support showed richer constructive participation, greater post-answer depth and, in coding-oriented conversations, greater persistence after observable breakdowns. These patterns indicate that scaffolding is part of the conversational environment in which constructive participation becomes more likely to appear. At the same time, scaffolded support was not randomly assigned. More difficult tasks, more motivated users or already constructive exchanges may elicit scaffolding while also producing constructive follow-up. The point is therefore not that exposure to scaffolding causes durable learning gains. Rather, the study identifies structured associations in naturalistic interaction and shows where they become visible: in conversations where support and user participation are mutually organized around continued work with ideas.

% The most actionable insight is that learning-oriented support is not a matter of adding more scaffolding in general. The broad label of scaffolding concealed important differences in support form. Feedback-like support and explanation were most clearly aligned with constructive engagement, whereas instructing and questioning were less aligned with visible elaboration. This distinction suggests that different support forms leave different roles for the user. Feedback responds to a user's prior attempt and can create an occasion for revision, evaluation or extension. Explanation can make mechanisms visible and give users material for testing their understanding. By contrast, detailed instructions can make the next step executable without requiring the user to articulate why it works, and questions in everyday assistant use may function as information gathering or clarification requests rather than as carefully designed prompts for reflection. This does not mean that instruction or questioning are unhelpful. They may be appropriate when users need direct task support, when information is missing or when the cost of continued exploration is high. But they illustrate why support quality for informal learning cannot be judged only by whether a response is pedagogical in form. The relevant question is whether the support preserves a productive role for the user in the work of understanding.

% The temporal analyses further sharpen this interpretation. Support and engagement were locally coupled rather than cumulatively amplified by default. Scaffolded turns were more often followed by constructive user turns in most settings, and users' current engagement state shaped the probability of receiving subsequent scaffolding. These patterns make everyday human--LLM interaction look less like a one-way instructional pipeline and more like a coupled sequence in which user state and assistant support condition one another. This temporal boundary is important. If scaffolded support simply produced cumulative learning-oriented engagement, one would expect constructive participation to rise steadily after the first scaffolded turn. Instead, longer-span contrasts were mixed. The strongest claim is therefore scale-specific: everyday LLM interaction repeatedly creates local moments in which support and user engagement align, but those moments do not automatically accumulate into a uniform upward trajectory. Informal learning with LLMs should be understood as a sequence of contingent opportunities, not as a guaranteed consequence of interacting with an intelligent assistant.

% These findings have design implications for personal AI assistants that aim to support lifelong learning in the flow of everyday work. The design target is not a generic tutor mode, and it is not simply more instructional language. A useful assistant must decide what kind of role to leave for the user at a given moment. Sometimes the right action is to answer directly, because the user primarily needs task completion or because a pedagogical detour would impose unnecessary friction. Sometimes the right action is to explain a mechanism, because the user is trying to understand why something works. Sometimes it is to give feedback on a user's attempt, because the user has already produced an idea that can be evaluated, revised or extended. Sometimes it is to offer a hint or partial cue, preserving the user's responsibility for the next step. And sometimes it may be to step back, because additional instruction would replace rather than support the user's own sense-making. This implies that future assistants need policies for calibrating support along at least three dimensions: the user's current framing, the task ecology and the user's local interaction state. A learning-oriented assistant should not merely be better at teaching; it should be better at recognizing when teaching is needed, what form it should take and how much cognitive responsibility should remain with the user.

% Several limits define the boundary of these claims. First, the analyses are observational, so scaffolded support should be interpreted as associated with engagement rather than causally responsible for it. Selection into support remains a central threat: complex tasks or motivated users may both elicit scaffolding and produce constructive follow-up. The adjacent-turn analyses reduce some ambiguity by conditioning on the immediately preceding interaction state, but they cannot eliminate unobserved time-varying confounding. Second, the labels capture behavioural manifestations of learning-oriented engagement rather than verified learning outcomes or durable knowledge change. Constructive engagement is a theoretically grounded behavioural proxy, not a direct measure of what users learned. Third, fine-grained support-form labels are not mutually exclusive and are noisier than the top-level distinction between scaffolded and non-scaffolded support. Fourth, the analysis focuses on coding and writing as two common everyday task ecologies; other domains may organize different opportunities for engagement and support. Finally, platform and model labels are not experimental assignments, so platform- or model-specific patterns should be treated as robustness evidence rather than causal comparisons. These limits define the appropriate interpretation: the study does not show that LLM support causes learning, but it shows how learning-relevant interaction can be observed, characterized and designed for in everyday AI use.

% Overall, the study suggests that LLMs are becoming part of the ordinary environments in which people encounter, test and extend understanding while working. They are not only instruments for faster task completion, and they are not educational technologies only when explicitly deployed as tutors. They are increasingly embedded in the everyday sequences through which people clarify problems, revise ideas and decide what to try next. Understanding this middle space is necessary for building future assistants that support lifelong learning without assuming that every user has entered a classroom. The central challenge is to make everyday human--AI collaboration a setting in which people can solve real problems while continuing to build, test and extend what they understand.

% Below is the previou section
%%%%%%%%%%%%%%%%%%%%%%%%%%%%%%%%%%%%%%%%%%%%%%%%%%%%%%
% The main finding is therefore not simply that learning-oriented engagement exists, but that constructive engagement is ecological. Everyday LLM use should not be treated as a single behavioural environment. Different task ecologies create different opportunities for users to ask questions, test assumptions, refine solutions and extend ideas. Coding-oriented interactions made uncertainty and failure more visible through errors, constraints and executable tests, whereas writing-oriented interactions more often involved drafting, editing or applying revisions without requiring users to articulate the underlying rationale. Similarly, longer exchanges should not be read as learning trajectories by themselves. They provide more opportunities for constructive turns and may reflect more complex tasks or more motivated users. The substantive point is that constructive engagement becomes most visible when everyday AI use shifts from answer delivery to sustained sense-making around the assistant's response.

% Assistant scaffolding was part of this interactional ecology, but its role should be interpreted carefully. Conversations containing scaffolded support showed higher constructive ratios, greater post-answer depth and, in coding-oriented conversations, greater persistence after observable breakdowns. Adjusted models indicated that these associations were not fully explained by measured differences in conversation composition. However, scaffolded support was not randomly assigned. More difficult tasks, more motivated users or already constructive exchanges may all elicit scaffolding while also producing constructive follow-up. The contribution is therefore not to show that scaffolding causes durable learning gains, but to identify structured associations in naturalistic use: scaffolded support marks conversational environments in which constructive participation is more likely to appear.

% A central implication is that learning-oriented support is not a matter of adding more scaffolding in general. The broad label of scaffolding concealed important differences in support form. Feedback-like support and explanation were most clearly aligned with constructive engagement, whereas instructing and questioning were less aligned with visible elaboration. This distinction matters because different support forms can preserve or reduce the user's opportunity to do cognitive work. Feedback can respond to a user's prior attempt and create an occasion for revision, evaluation or extension. Explanation can make mechanisms visible and support users in testing their understanding. By contrast, detailed instructions may make the next step executable without requiring the user to articulate why it works, and questions in everyday assistant use may function as clarification requests rather than designed prompts for reflection. Future assistants that aim to support learning should therefore calibrate not only whether to provide support, but what form of support to provide given the user's framing and the task ecology.

% The temporal analyses further suggest that informal learning with LLMs is best understood as a local interaction process. Scaffolded turns were more often followed by constructive user turns in most settings, and users' current engagement state shaped the probability of receiving subsequent scaffolding. Support and engagement therefore appeared as coupled states, not as independent events in a one-way instructional pipeline. This finding also guards against over-reading conversation-level associations. Longer-span before--after contrasts around the first scaffolded turn were mixed, arguing against a simple cumulative-uplift account in which support steadily raises constructive engagement over the remainder of a conversation. The strongest temporal claim is scale-specific: everyday LLM interaction repeatedly creates local opportunities for constructive engagement, but those opportunities do not automatically accumulate into a uniform upward trajectory.

% These results have implications for the design of personal AI assistants. Systems that support informal learning should not simply turn every interaction into a lesson or add more instructional language. They should recognize when a user primarily needs direct assistance, when a user is active but not yet constructively engaged, and when a user is already working with ideas and may benefit from feedback or explanation that sustains elaboration. This requires domain-sensitive and state-sensitive support policies. In some moments, the best support may be an explanation that helps users understand a mechanism; in others, it may be feedback on a user's attempt, a partial cue that keeps responsibility with the user or a direct answer that avoids an unnecessary pedagogical detour. The design target is not a generic tutor mode, but adaptive support that helps people build knowledge while solving real problems.

% Several limits define the boundary of these claims. First, the analyses are observational, so scaffolding should be interpreted as associated with engagement rather than causally responsible for it. Second, the labels capture behavioural manifestations of learning-oriented engagement rather than verified learning outcomes or durable knowledge change. Constructive engagement is a theoretically grounded behavioural proxy, not a direct measure of what users learned. Third, fine-grained support-form labels are not mutually exclusive and are noisier than the top-level distinction between scaffolded and non-scaffolded support. Fourth, the analysis focuses on coding and writing as two common everyday task ecologies; other domains may organize different opportunities for engagement and support. Finally, model labels and platform contexts are not experimental assignments, so model- or platform-specific comparisons should be treated as robustness checks rather than causal evidence. These limits are important, but they do not erase the central contribution: the study makes informal learning in everyday AI use observable at scale and shows that learning-oriented engagement depends on the situated alignment of task ecology, user framing, support form and interaction state.

% Overall, the study suggests that LLMs are becoming part of the ordinary environments in which people encounter, test and extend understanding while working. They are not only tools for faster task completion, and they are not educational technologies only when explicitly deployed as tutors. They are increasingly embedded in the everyday sequences through which people clarify problems, revise ideas and decide what to try next. Understanding this middle space is necessary for designing future assistants that support lifelong learning without assuming that every user has entered a classroom.